\section{Introduction}

Determining whether one equational law implies another is a basic operation in universal algebra and automated theorem proving. In the Equational Theories Project (ETP), this question appears at scale: a large implication graph over magma laws has been computed and formalized, making implication prediction and certification a central practical task~\cite{bolan2025etp}. Fast implication decisions are useful as a routing layer before expensive proof search.

Classical finite-model tools such as Mace4 emphasize falsification by small counterexamples~\cite{mccune2003mace4}. In parallel, learning-based approaches report strong predictive performance for equational tasks~\cite{piepenbrock2021learning,axelrod2026twitch}, and latent-geometry analyses reveal structure in the space of equational theories~\cite{berlioz2026latent}. For Stage-1 settings, however, there remains a need for a transparent, reproducible baseline that always emits strict boolean answers while keeping runtime low.

We study single-premise implication queries over magmas: given equations $E_1$ and $E_2$, decide whether $E_1 \models E_2$ with output in $\{\texttt{true},\texttt{false}\}$. Our solver uses a countermodel-first hybrid policy: bounded finite model search up to order $3$, then bounded symbolic derivation, and finally a deterministic boolean fallback required by the task contract.

Our contributions are threefold.
\begin{enumerate}[leftmargin=*]
\item We formalize the Stage-1 policy and prove structural guarantees: totality, strict boolean output, and soundness of negative answers.
\item We provide reproducible computational evidence on a curated benchmark of ten implication pairs, obtaining $10/10$ accuracy with mean runtime $0.0623$ seconds per query.
\item We identify a concrete limitation of the current derivation module: one true case is solved only by fallback, which motivates stronger positive-proof mechanisms.
\end{enumerate}

Section~2 introduces notation and the policy definition. Section~3 states and proves the main theoretical and computational results. Section~4 discusses implications, limitations, and open questions. Section~5 concludes with directions for extended evaluation and prover integration.
